\subsection{Separating screening width from the voting rule}
The first forest comparison changed both the feature-selection rule and the number of columns available to the forest. We followed it with four otherwise identical pipelines. One fitted the forest selector to all columns. The other three first kept 40\% by pivotal minPI, by $F$ ranking, or by uniform random sampling. Every forest then returned 10\% of the original width to the same downstream learner. The eight previously examined benchmark sources form a retrospective comparison. Two additional sources were specified before fitting any of the four pipelines to them.

For Gisette and ISOLET, all three ways of selecting the first 40\% met the 20\% time criterion and the quality tolerances against the direct forest on all five splits for each downstream learner. On Gisette, pivotal-to-$F$ time ratios were 0.937 for ridge and 0.953 for trees. Relative to random screening they were 1.059 and 1.010. By contrast, the pivotal-to-unscreened ratios in these fresh fits were 0.415 and 0.468. Across all ten sources, neither learner produced a seed for which the pivotal screen followed by ET was at least 20\% faster than the equal-width $F$ screen followed by ET while meeting the quality tolerances against that comparator (0 of 100 source--learner--seed pairs). The timing gains against direct ET are consistent with reducing the input width of the forest.

CNAE-9 contains 1,080 documents with 856 word-frequency features and nine classes~\citep{ciarelli2009cnae}. On five stratified row splits, voting followed by forest selection reached mean accuracy 0.874 with ridge and 0.861 with trees. With $F$ screening these values were 0.918 and 0.897, while random screening reached 0.710 and 0.717. The choice of the first-stage columns affected both learners on this source. Of the 270 held-out documents per split, 12--16 had a word-frequency vector identical to one in training, so these results include repeated representations. Removing these test duplicates without refitting left the pivotal-minus-$F$ AUC difference negative on every split with both learners (mean differences $-0.010$ and $-0.012$). SECOM contains semiconductor manufacturing records with sensor measurements, missing values and process pass/fail labels~\citep{mccann2008secom}. We split by acquisition time and fitted median imputation on training records only. Among 392 later records, 24 were process failures. With a final tree model, voting gave mean AUC 0.531 and average precision 0.075, compared with 0.440 and 0.058 after $F$ screening, 0.466 and 0.061 after random screening, and 0.504 and 0.065 without screening. Ridge reversed the AUC comparison with $F$ (0.549 versus 0.564). Default-threshold balanced accuracy was close to 0.5. The five seeds reuse these same 24 positive test records; their mean describes the fitted procedures on a single temporal cohort.
